\begin{theorem}[machine-verified]\label{thm:orders}
For every even $k$ listed in Table~\ref{tab:orders}, no loopless cubic graph on $k$ vertices colors
the corresponding graph with a vertex map of kind \textup{(O)}, \textup{(E0)} or \textup{(E1)}.
The list contains every even $k$ with $40 \le k \le 50$ for both $\Ga$ and $\Gb$%%EXTRA_ORDERS%%.
\end{theorem}

Goedgebeur et al.\ established by an exhaustive computation that every bridgeless cubic graph on
at most 38 vertices has a Petersen coloring~\cite[Observation~9]{GJMMMU2026}. The statement
extends to graphs with parallel edges.

\begin{lemma}\label{lem:digon}
If a connected bridgeless cubic graph with parallel edges has no Petersen coloring, then a
connected bridgeless cubic graph of smaller order has none either.
\end{lemma}

\begin{proof}
Let $x, y$ be joined by two parallel edges; they are not joined by three, since the graph on two
vertices is Petersen colorable. Let $xx'$ and $yy'$ be the remaining edges at $x$ and $y$. Then
$x' \ne y'$, because otherwise the third edge at $x'$ would be a bridge. Delete $x$ and $y$ and
add the edge $x'y'$. Every edge cut of the new graph corresponds to an edge cut of the same size
of the old one, so the new graph is connected, bridgeless and cubic. If it had a Petersen
coloring $\sigma$ with $\sigma(x'y') = m$, then labelling $xx'$ and $yy'$ by $m$ and the two
parallel edges by the two other edges at one end of $m$ would give a Petersen coloring of the
old graph.
\end{proof}

\begin{theorem}\label{thm:h3}
The graphs $\Ga$ and $\Gb$ are colorable only by themselves: both belong to $\Hthree$.
\end{theorem}

\begin{proof}
Let $G$ be one of the two graphs and suppose that some connected bridgeless cubic graph of order
less than 52 colors $G$. By Corollary~\ref{cor:modes} there is one, say $H$, of even order
$k < 52$ with a vertex map of kind (O), (E0) or (E1). Since $H$ colors a graph without a Petersen
coloring, $H$ has no Petersen coloring. By Lemma~\ref{lem:digon}, applied repeatedly, there is a
simple connected bridgeless cubic graph of order at most $k$ without a Petersen coloring, so
$k \ge 40$ by~\cite[Observation~9]{GJMMMU2026}. This contradicts Theorem~\ref{thm:orders}.
Theorem~\ref{thm:mmsw} gives the conclusion.
\end{proof}

Theorem~\ref{thm:orders} does not depend on the computation of~\cite{GJMMMU2026}; only the orders
below 40 in Theorem~\ref{thm:h3} do. The orders below 40 that are listed in
Table~\ref{tab:orders} were refuted directly, and the remaining ones were not decided within six
hours each.
